\chapter{Finite Volumes: Numerical Fluxes and High-Resolution Schemes}\label{ch:8}
\begin{chaptersummary}
The weak form of a conservation law describes discontinuities through control-volume flux balances. Cell averages, Riemann problems, and shared fluxes lead to finite-volume methods; Godunov's order barrier motivates nonlinear reconstruction. Conservation, monotonicity, and entropy stability impose distinct constraints on fluxes and updates, providing the basis for DG and FR interface design.
\end{chaptersummary}
\learninggoals{Construct cell updates from integral conservation; explain Riemann fluxes; distinguish Godunov's theorem, TVD, and high-order reconstruction; check reconstructed averages and units.}
\studyroute{The chapter develops Godunov, MUSCL, TVD, and WENO in sequence. Characteristic decomposition for systems, full smoothness indicators, and nonlinear entropy-convergence proofs are left for further study.}{finite volume; Riemann problem; Godunov order barrier; MUSCL; WENO}{Calculate a minmod reconstruction from three cell averages, then apply the Rusanov flux to two adjacent control volumes and check the signs of their internal fluxes.}
\section{Control Volumes, Cell Averages, and Shared Fluxes}\index{Finite volume method}
Finite-volume methods start from control-volume conservation:
\[
\frac{\dd}{\dd t}
\int_{\Omega_i}U\dd V
+
\int_{\partial\Omega_i}
\bm F\cdot\bm n\dd S
=
\int_{\Omega_i}S\dd V.
\]
Adjacent cells use the same interface flux with opposite signs. Internal fluxes therefore cancel in a global sum, giving local and global conservation.


Define $\bar U_K=|K|^{-1}\int_KU$. Degrees of freedom are cell averages; interface states are reconstructed from them.
A consistent numerical normal flux satisfies $\widehat F(U,U;n)=F(U)\cdot n$ and changes sign when the states are exchanged and the normal is reversed.
For a one-dimensional scalar equation, Riemann initial data are constant states $U_L,U_R$ on either side of an interface.
The exact interface flux is the Godunov flux; approximate solvers use computable wave-speed estimates.
For example,
\[
\widehat F(U_L,U_R)=\frac{F(U_L)+F(U_R)}2-\frac{\alpha}2(U_R-U_L),
\qquad
\alpha\ge\max_{[U_L,U_R]}|F'(U)|.
\]
For linear advection $F(U)=aU$ with $a>0$, the exact flux is $aU_L$.
A cell update also integrates the flux over one time step. High-order spatial reconstruction and time integration jointly determine fully discrete accuracy.

\section{Godunov's Theorem and Nonlinear Reconstruction}
Godunov's theorem is a basic limitation in numerical conservation laws. For linear scalar advection, consider a scheme that is

\begin{enumerate}
  \item linear;
  \item monotone, creating no new extrema;
  \item consistent.
\end{enumerate}

For general noninteger Courant numbers, a linear monotone consistent scheme has at most first-order accuracy. Exact grid translation at integer Courant numbers for constant-speed advection is a special case.

Higher-order linear reconstruction generally requires alternating positive and negative weights, whereas monotonicity requires a convex-combination update. These requirements are generally incompatible.

Improving the accuracy of monotone schemes therefore requires changing these conditions. Modern high-resolution methods commonly use nonlinear reconstruction:

\begin{authornote}
\textbf{Potential Question:}
"Why does Godunov's theorem motivate nonlinear reconstruction? Which restriction does nonlinearity remove?"

\textbf{Author Supplement:}
Under the theorem's assumptions, linear monotone schemes are at most first order. High-resolution methods therefore make reconstruction or limiting depend on the local solution: retain higher order in smooth regions and adjust slopes near steep gradients to suppress oscillation. The scheme becomes nonlinear and leaves the class of linear schemes restricted by the theorem.

\textbf{Core Conceptual Connection:}
\[
\text{Godunov theorem}
\longrightarrow
\text{nonlinear reconstruction}
\longrightarrow
\text{high-resolution schemes}
\]
\end{authornote}

\begin{itemize}
  \item Preserve conservation.
  \item Allow a nonlinear algorithm.
  \item Recover high order in smooth regions.
  \item Reduce to a first-order monotone scheme near discontinuities.
\end{itemize}

For a finite-volume method,
\[
\frac{\dd \bar U_i}{\dd t}
=
-\frac{1}{\Delta x}
\left(
\widehat F_{i+1/2}
-
\widehat F_{i-1/2}
\right).
\]
the key step is constructing left and right interface states
\[
U_{i+1/2}^{L},
\qquad
U_{i+1/2}^{R},
\]
and evaluating a numerical flux with an exact or approximate Riemann solver.

MUSCL uses piecewise-linear reconstruction, for example,
\[
U_{i+1/2}^{L}
=
U_i+\frac12\sigma_i,
\]

\begin{askai}
Ask AI:
\begin{quote}
Construct MUSCL reconstruction from three neighboring cell averages. Calculate \(U_{i+1/2}^L\) using an unlimited slope and a minmod limiter, and explain how limiting retains second order in smooth regions and reverts to first order near discontinuities.
\end{quote}

Follow-up:
\begin{quote}
For data containing a discontinuity, calculate and compare minmod, van Leer, and Superbee slopes.
\end{quote}
\end{askai}

Here $\sigma_i$ is the increment across one cell width and has the same units as $U_i$.
For example, $\sigma_i=\operatorname{minmod}(U_i-U_{i-1},U_{i+1}-U_i)$,
where minmod returns the argument of smaller magnitude for equal signs, and zero otherwise.
The slope \(\sigma_i\) adapts to local smoothness; minmod, van Leer, and Superbee impose different restrictions. This retains approximately second-order behavior in smooth regions and reduces reconstruction order near shocks or discontinuities to suppress oscillations.

A TVD method requires nonincreasing discrete total variation:
\[
TV(U^{n+1})\le TV(U^n),
\qquad
TV(U)=\sum_i|U_{i+1}-U_i|.
\]
This turns suppression of spurious oscillations into an analyzable discrete condition.

ENO/WENO selects or weights several candidate stencils rather than directly limiting one slope. WENO uses a high-order combination in smooth regions and reduces weights on stencils containing a discontinuity.

Godunov's theorem exposes a structural tradeoff:
\[
\text{High-order accuracy}
\quad\leftrightarrow\quad
\text{Monotonicity and nonoscillation}.
\]
Scheme design adjusts reconstruction and dissipation to local solution smoothness.

\section{Conservation, Entropy, and High Order in Smooth Regions}
WENO combines candidate reconstructions $q_r$ with nonlinear weights:
\[
U_{i+1/2}^{L}=\sum_r\omega_rq_r,\qquad
\omega_r=\frac{\alpha_r}{\sum_s\alpha_s},\qquad
\alpha_r=\frac{d_r}{(\epsilon_W+\beta_r)^2}.
\]
$\beta_r$ measures stencil smoothness and $d_r$ gives the ideal weights of the corresponding high-order linear combination;
the scale of $\epsilon_W$ must be compatible with the variables and mesh.
Smooth-region order depends on how rapidly the weights recover their ideal values; extrema require separate analysis.
Fully discrete scalar TVD results depend on the flux, limiter, and time step.
Entropy-stable systems additionally require entropy variables and compatible numerical entropy fluxes. For the original WENO construction, see~\cite{jiang1996}.
